\begin{thebibliography}{Klo86b}\small
\bibitem[Ada62]{Ada62} J. F. Adams, \emph{Vector fields on spheres}, Ann. of Math. (2) \textbf{75} (1962), 603--632. \url{https://www.math.ens.psl.eu/~benoist/refs/Adams.pdf}.
\bibitem[AP05]{AP05} A.-M.~Aubert and R.~Plymen, Plancherel measure for $\GL(n,F)$ and $\GL(m,D)$: explicit formulas and Bernstein decomposition, \emph{J.\ Number Theory} 112 (2005), 26--66.
\bibitem[Art83]{Art83} J.~Arthur, A Paley--Wiener theorem for real reductive groups, \emph{Acta Math.} 150 (1983), 1--89.
\bibitem[AT68]{AT68} E. Artin and J. Tate, \emph{Class Field Theory}, W.~A.~Benjamin, New York, 1968; revised edition, AMS Chelsea, Providence, 2009.
\bibitem[Ati66]{Ati66} M. F. Atiyah, \emph{Power operations in $K$-theory}, Quart. J. Math. Oxford (2) \textbf{17} (1966), 165--193. \url{https://doi.org/10.1093/qmath/17.1.165}.
\bibitem[Bad08]{Bad08} A.~I.~Badulescu, Global Jacquet--Langlands correspondence, multiplicity one and classification of automorphic representations, with an appendix by N.~Grbac, \emph{Invent.\ Math.} 172 (2008), 383--438; arXiv:0704.2920.
\bibitem[Bad02]{Bad02} A.~I.~Badulescu, Correspondance de Jacquet--Langlands pour les corps locaux de caract\'eristique non nulle, \emph{Ann.\ Sci.\ \'Ec.\ Norm.\ Sup.} 35 (2002), 695--747.
\bibitem[BD01]{BD01} J.-L.~Brylinski and P.~Deligne, Central extensions of reductive groups by $\mathbf K_2$, \emph{Publ.\ Math.\ IH\'ES} 94 (2001), 5--85.
\bibitem[Ber84]{Ber84} J.~N.~Bernstein, Le ``centre'' de Bernstein (r\'edig\'e par P.~Deligne), in \emph{Repr\'esentations des groupes r\'eductifs sur un corps local}, Hermann, 1984, 1--32.
\bibitem[BH06]{BH06} C. J. Bushnell and G. Henniart, \emph{The local Langlands conjecture for $\mathrm{GL}(2)$}, Grundlehren 335, Springer, 2006.
\bibitem[BH11]{BH11} C. J. Bushnell and G. Henniart, The essentially tame Jacquet--Langlands correspondence for inner forms of $\mathrm{GL}(n)$, \emph{Pure Appl. Math. Q.} \textbf{7} (2011), 469--538.
\bibitem[BLS99]{BLS99} W.~D.~Banks, J.~Levy and M.~R.~Sepanski, Block-compatible metaplectic cocycles, \emph{J.\ reine angew.\ Math.} 507 (1999), 131--163.
\bibitem[BS08]{BS08} M. Boyarchenko and M. Sabitova, The orbit method for profinite groups and a $p$-adic analogue of Brown's theorem, \emph{Israel J. Math.} \textbf{165} (2008), 67--91.
\bibitem[BZ76]{BZ76} I.~N.~Bernstein and A.~V.~Zelevinsky, Representations of the group $\GL(n,F)$ where $F$ is a local non-archimedean field, \emph{Russian Math.\ Surveys} 31:3 (1976), 1--68.
\bibitem[BZ77]{BZ77} I.~N.~Bernstein and A.~V.~Zelevinsky, Induced representations of reductive $p$-adic groups.\ I, \emph{Ann.\ Sci.\ \'Ec.\ Norm.\ Sup.} 10 (1977), 441--472.
\bibitem[CD84]{CD84} L. Clozel and P. Delorme, \emph{Le th\'eor\`eme de Paley--Wiener invariant pour les groupes de Lie r\'eductifs}, Invent. Math. \textbf{77} (1984), 427--453.
\bibitem[Coh18]{Coh18} J. Cohen, \emph{Transfer of representations and orbital integrals for inner forms of $\GL_n$}, Canad. J. Math. \textbf{70} (2018), 595--627. \url{https://arxiv.org/abs/1606.02141}.
\bibitem[Cor74]{Cor74} L. Corwin, Representations of division algebras over local fields, \emph{Adv. Math.} \textbf{13} (1974), 259--267.
\bibitem[DKV84]{DKV84} P. Deligne, D. Kazhdan, and M.-F. Vign\'eras, Repr\'esentations des alg\`ebres centrales simples $p$-adiques, in: \emph{Repr\'esentations des groupes r\'eductifs sur un corps local}, Hermann, Paris, 1984, 33--117.
\bibitem[DM78]{DM78} J.~Dixmier and P.~Malliavin, Factorisations de fonctions et de vecteurs ind\'efiniment diff\'erentiables, \emph{Bull.\ Sci.\ Math.} 102 (1978), 307--330.
\bibitem[FK86]{FK86} Y.~Z.~Flicker and D.~A.~Kazhdan, Metaplectic correspondence, \emph{Publ.\ Math.\ IH\'ES} 64 (1986), 53--110.
\bibitem[FK88]{FK88} Y.~Z.~Flicker and D.~A.~Kazhdan, A simple trace formula, \emph{J.\ Analyse Math.} 50 (1988), 189--200.
\bibitem[Fli80]{Fli80} Y.~Z.~Flicker, Automorphic forms on covering groups of $\GL(2)$, \emph{Invent.\ Math.} 57 (1980), 119--182.
\bibitem[FV02]{FV02} I. B. Fesenko and S. V. Vostokov, \emph{Local fields and their extensions}, 2nd ed., Transl. Math. Monogr. 121, AMS, 2002.
\bibitem[Gao12]{Gao12} F.~Gao, The residual spectrum of $\mathrm{Mp}_4(\A_k)$, arXiv:1212.0617.
\bibitem[Gao18]{Gao18} F. Gao, The Langlands--Shahidi $L$-functions for Brylinski--Deligne extensions, \emph{Amer. J. Math.} \textbf{140} (2018), 83--137.
\bibitem[Gel76]{Gel76} S.~S.~Gelbart, \emph{Weil's Representation and the Spectrum of the Metaplectic Group}, Lecture Notes in Math.\ 530, Springer, 1976.
\bibitem[Ger77]{Ger77} P. G\'erardin, Weil representations associated to finite fields, \emph{J. Algebra} \textbf{46} (1977), 54--101.
\bibitem[GG18]{GG18} W.~T.~Gan and F.~Gao, The Langlands--Weissman program for Brylinski--Deligne extensions, \emph{Ast\'erisque} 398 (2018), 187--275.
\bibitem[GGPS69]{GGPS69} I.~M.~Gelfand, M.~I.~Graev and I.~I.~Piatetski-Shapiro, \emph{Representation Theory and Automorphic Functions}, Saunders, 1969.
\bibitem[GGW18]{GGW18} W.~T.~Gan, F.~Gao and M.~H.~Weissman, $L$-groups and the Langlands program for covering groups: a historical introduction, \emph{Ast\'erisque} 398 (2018), 1--31.
\bibitem[God66]{God66} R.~Godement, The spectral decomposition of cusp-forms, in \emph{Algebraic Groups and Discontinuous Subgroups}, Proc.\ Sympos.\ Pure Math.\ 9, Amer.\ Math.\ Soc., 1966, 225--234.
\bibitem[GT16]{GT16} W. T. Gan and S. Takeda, A proof of the Howe duality conjecture, \emph{J. Amer. Math. Soc.} \textbf{29} (2016), 473--493.
\bibitem[HC54]{HC54} Harish-Chandra, Representations of semisimple Lie groups.\ II, \emph{Trans.\ Amer.\ Math.\ Soc.} 76 (1954), 26--65.
\bibitem[HC68]{HC68} Harish-Chandra, \emph{Automorphic Forms on Semisimple Lie Groups}, Lecture Notes in Math.\ 62, Springer, 1968.
\bibitem[HC70]{HC70} Harish-Chandra, \emph{Harmonic Analysis on Reductive $p$-adic Groups} (notes by G.~van Dijk), Lecture Notes in Math.\ 162, Springer, 1970.
\bibitem[HC78]{HC78} Harish-Chandra, Admissible invariant distributions on reductive $p$-adic groups, \emph{Queen's Papers in Pure and Appl.\ Math.} 48 (1978), 281--347.
\bibitem[Hil08]{Hil08} R.~Hill, Geometric construction of metaplectic covers of $\GL_n$ in characteristic zero, arXiv:0708.0108.
\bibitem[How74]{How74} R.~Howe, The Fourier transform and germs of characters (case of $\GL_n$ over a $p$-adic field), \emph{Math.\ Ann.} 208 (1974), 305--322.
\bibitem[How77a]{How77a} R. Howe, Kirillov theory for compact $p$-adic groups, \emph{Pacific J. Math.} \textbf{73} (1977), 365--381.
\bibitem[How77b]{How77b} R. Howe, Tamely ramified supercuspidal representations of $\mathrm{GL}_n$, \emph{Pacific J. Math.} \textbf{73} (1977), 437--460.
\bibitem[How79]{How79} R. Howe, $\theta$-series and invariant theory, in: \emph{Automorphic forms, representations and $L$-functions}, Proc. Sympos. Pure Math. 33, Part 1, AMS, 1979, 275--285.
\bibitem[Isa76]{Isa76} I. M. Isaacs, \emph{Character theory of finite groups}, Academic Press (1976), Ch.~8 (Brauer's characterization of characters).
\bibitem[J88]{J88} D.~Joyner, On the metaplectic analog of Kazhdan's ``endoscopic'' lifting, \emph{Israel J.\ Math.} 61 (1988), 113--154.
\bibitem[Jai06]{Jai06} A. Jaikin-Zapirain, Zeta function of representations of compact $p$-adic analytic groups, \emph{J. Amer. Math. Soc.} \textbf{19} (2006), 91--118.
\bibitem[JC26]{JC26} D.~Joyner and Claude, On the metaplectic analog of Kazhdan's ``endoscopic'' lifting (original 1988, corrected 2026), \texttt{metaplectic\_kazhdan\_revised12.tex}.
\bibitem[JL70]{JL70} H.~Jacquet and R.~P.~Langlands, \emph{Automorphic Forms on $\GL(2)$}, Lecture Notes in Math.\ 114, Springer, 1970.
\bibitem[Kaz86]{Kaz86} D.~Kazhdan, Cuspidal geometry of $p$-adic groups, \emph{J.\ Analyse Math.} 47 (1986), 1--36.
\bibitem[Kaz84]{Kaz84} D.~Kazhdan, On lifting, in \emph{Lie Group Representations II}, Lecture Notes in Math.\ 1041, Springer, 1984, 209--249.
\bibitem[Klo86a]{Klo86a} J.~Klose, Kommutatoren und 2-Cohomologie $p$-adischer Quaternionenschiefk\"orper, \emph{Math.\ Z.} 191 (1986), 261--282.
\bibitem[Klo86b]{Klo86b} J.~Klose, Metaplektische Erweiterungen von Quaternionenschiefk\"orpern, \emph{Math.\ Z.} 193 (1986), 625--649.
\bibitem[Kot05]{Kot05} R.~E.~Kottwitz, Harmonic analysis on reductive $p$-adic groups and Lie algebras, in \emph{Harmonic Analysis, the Trace Formula, and Shimura Varieties}, Clay Math.\ Proc.\ 4, Amer.\ Math.\ Soc., 2005, 393--522.
\bibitem[Kot84]{Kot84} R. E. Kottwitz, \emph{Stable trace formula: cuspidal tempered terms}, Duke Math. J. \textbf{51} (1984), no.~3, 611--650. \url{https://doi.org/10.1215/S0012-7094-84-05129-9}.
\bibitem[Kot88]{Kot88} R.~E.~Kottwitz, Tamagawa numbers, \emph{Ann.\ of Math.} 127 (1988), 629--646.
\bibitem[KP84]{KP84} D.~Kazhdan and S.~J.~Patterson, Metaplectic forms, \emph{Publ.\ Math.\ IH\'ES} 59 (1984), 35--142; Corrections, \emph{ibid.} 62 (1985), 203.
\bibitem[KP86]{KP86} D.~Kazhdan and S.~J.~Patterson, Towards a generalized Shimura correspondence, \emph{Adv.\ in Math.} 60 (1986), 161--234.
\bibitem[KT26]{KT26} O.~Kivinen and C.-C.~Tsai, Shalika germs for tamely ramified elements in $\GL_n$, arXiv:2209.02509v4 (August 13, 2026). (In an earlier version, parabolic descent was Proposition~2.42.)
\bibitem[Kub67]{Kub67} T.~Kubota, Topological covering of $\SL(2)$ over a local field, \emph{J.\ Math.\ Soc.\ Japan} 19 (1967), 114--121.
\bibitem[Kub69]{Kub69} T.~Kubota, \emph{On Automorphic Functions and the Reciprocity Law in a Number Field}, Lectures in Math., Kyoto Univ.\ 2, Kinokuniya, Tokyo, 1969.
\bibitem[KZ80]{KZ80} H. Koch and E.-W. Zink, Zur Korrespondenz von Darstellungen der Galoisgruppen und der zentralen Divisionsalgebren \"uber lokalen K\"orpern (der zahme Fall), \emph{Math. Nachr.} \textbf{98} (1980), 83--119.
\bibitem[Lan02]{Lan02} S.~Lang, \emph{Algebra}, revised 3rd ed., Graduate Texts in Math.\ 211, Springer, 2002.
\bibitem[Laz65]{Laz65} M. Lazard, \emph{Groupes analytiques $p$-adiques}, Publ. Math. IH\'ES \textbf{26} (1965); see also J. D. Dixon, M. P. F. du Sautoy, A. Mann and D. Segal, \emph{Analytic pro-$p$ groups}, 2nd ed., Cambridge Univ. Press (1999), Thm.~5.1.5 and Ch.~4.
\bibitem[Li10]{Li10} W.-W.~Li, La formule des traces pour les rev\^etements de groupes r\'eductifs connexes.\ I: Le d\'eveloppement g\'eom\'etrique fin, arXiv:1004.4011.
\bibitem[Li12]{Li12} W.-W.~Li, La formule des traces pour les rev\^etements de groupes r\'eductifs connexes.\ II: Analyse harmonique locale, \emph{Ann.\ Sci.\ \'Ec.\ Norm.\ Sup\'er.} (4) 45 (2012), 787--859.
\bibitem[Li14]{Li14} W.-W.~Li, La formule des traces pour les rev\^etements de groupes r\'eductifs connexes.\ IV: Distributions invariantes, \emph{Ann.\ Inst.\ Fourier} 64 (2014), 2379--2448.
\bibitem[Loe08]{Loe08} F. Loeser, \emph{Seattle lectures on motivic integration}, Proposition~1.2.1 (the $p$-adic change-of-variables formula), \url{https://webusers.imj-prg.fr/~francois.loeser/notes_seattle_09_04_2008.pdf}.
\bibitem[Mat69]{Mat69} H.~Matsumoto, Sur les sous-groupes arithm\'etiques des groupes semi-simples d\'eploy\'es, \emph{Ann.\ Sci.\ \'Ec.\ Norm.\ Sup.} 2 (1969), 1--62.
\bibitem[McN12]{McN12} P. J. McNamara, Principal series representations of metaplectic groups over local fields, in: \emph{Multiple Dirichlet series, L-functions and automorphic forms}, Progr. Math. 300, Birkh\"auser, 2012, 299--327.
\bibitem[Mez01]{Mez01} P.~Mezo, Comparisons of general linear groups and their metaplectic coverings.\ II, \emph{Represent.\ Theory} 5 (2001), 524--580.
\bibitem[Mez02]{Mez02} P.~Mezo, Comparisons of general linear groups and their metaplectic coverings.\ I, \emph{Canad.\ J.\ Math.} 54 (2002), 92--137.
\bibitem[Mez04]{Mez04} P.~Mezo, Metaplectic tensor products for irreducible representations, \emph{Pacific J.\ Math.} 215 (2004), no.~1.
\bibitem[Mil71]{Mil71} J.~Milnor, \emph{Introduction to Algebraic $K$-Theory}, Ann.\ of Math.\ Studies 72, Princeton Univ.\ Press, 1971.
\bibitem[MilADT06]{MilADT06} J. S. Milne, \emph{Arithmetic Duality Theorems}, 2nd ed., BookSurge, 2006, Chapter~I, \S4. \url{https://www.jmilne.org/math/Books/ADTnot.pdf}.
\bibitem[Moo68]{Moo68} C.~C.~Moore, Group extensions of $p$-adic and adelic linear groups, \emph{Publ.\ Math.\ IH\'ES} 35 (1968), 5--70.
\bibitem[MVW87]{MVW87} C. M\oe glin, M.-F. Vign\'eras, and J.-L. Waldspurger, \emph{Correspondances de Howe sur un corps $p$-adique}, Lecture Notes in Math. 1291, Springer, 1987.
\bibitem[MW95]{MW95} C.~M\oe glin and J.-L.~Waldspurger, \emph{Spectral Decomposition and Eisenstein Series}, Cambridge Tracts in Math.\ 113, Cambridge Univ.\ Press, 1995.
\bibitem[Neu99]{Neu99} J.~Neukirch, \emph{Algebraic Number Theory}, Grundlehren 322, Springer, 1999.
\bibitem[Pat14]{Pat14} S.~P.~Patel, A question on splitting of metaplectic covers, arXiv:1406.3978.
\bibitem[PRa88]{PRa88} G.~Prasad and M.~S.~Raghunathan, Topological central extensions of $\mathrm{SL}_1(D)$, \emph{Invent.\ Math.} 92 (1988), 645--689.
\bibitem[PRp96]{PRp96} G.~Prasad and A.~S.~Rapinchuk, Computation of the metaplectic kernel, \emph{Publ.\ Math.\ IH\'ES} 84 (1996), 91--187.
\bibitem[Rap86]{Rap86} A.~S.~Rapinchuk, The metaplectic kernel for the group $\mathrm{SL}(1,D)$ (Russian), \emph{Dokl.\ Akad.\ Nauk BSSR} 30 (1986), 197--200.
\bibitem[Rap88]{Rap88} A.~S.~Rapinchuk, Multiplicative arithmetic of division algebras over number fields and the metaplectic problem, \emph{Math.\ USSR Izv.} 31 (1988), 349--379.
\bibitem[Rei75]{Rei75} I. Reiner, \emph{Maximal orders}, Academic Press, 1975.
\bibitem[Rie70]{Rie70} C. Riehm, \emph{The norm $1$ group of a $\mathfrak p$-adic division algebra}, Amer. J. Math. \textbf{92} (1970), 499--523.
\bibitem[Rog81]{Rog81} J.~D.~Rogawski, An application of the building to orbital integrals, \emph{Compositio Math.} 42 (1981), 417--423.
\bibitem[Rog83]{Rog83} J.~D.~Rogawski, Representations of $\GL(n)$ and division algebras over a $p$-adic field, \emph{Duke Math.\ J.} 50 (1983), 161--196.
\bibitem[Sav88]{Sav88} G.~Savin, Local Shimura correspondence, \emph{Math.\ Ann.} 280 (1988), 185--190.
\bibitem[Seg68]{Seg68} G. Segal, \emph{Equivariant $K$-theory}, Publ. Math. IH\'ES \textbf{34} (1968), 129--151. \url{https://doi.org/10.1007/BF02684593}.
\bibitem[Ser79]{Ser79} J.-P.~Serre, \emph{Local Fields}, Graduate Texts in Math.\ 67, Springer, 1979.
\bibitem[Sha72]{Sha72} J.~A.~Shalika, A theorem on semi-simple $\mathcal P$-adic groups, \emph{Ann.\ of Math.} 95 (1972), 226--242.
\bibitem[Shi73]{Shi73} G.~Shimura, On modular forms of half integral weight, \emph{Ann.\ of Math.} 97 (1973), 440--481.
\bibitem[Sur05]{Sur05} B.~Sury, Central extensions of a $p$-adic division algebra, in \emph{Number Theory}, Ramanujan Math.\ Soc.\ Lect.\ Notes Ser.\ 1 (2005), 143--153.
\bibitem[SW00]{SW00} P. Symonds and T. Weigel, Cohomology of $p$-adic analytic groups, in: \emph{New horizons in pro-$p$ groups}, Progr. Math. 184, Birkh\"auser, 2000, 349--410.
\bibitem[vD72]{vD72} G.~van Dijk, Computation of certain induced characters of $p$-adic groups, \emph{Math.\ Ann.} 199 (1972), 229--240.
\bibitem[Wal88]{Wal88} N.~R.~Wallach, \emph{Real Reductive Groups I}, Pure and Applied Math.\ 132, Academic Press, 1988.
\bibitem[Wal03]{Wal03} J.-L.~Waldspurger, La formule de Plancherel pour les groupes $p$-adiques (d'apr\`es Harish-Chandra), \emph{J.\ Inst.\ Math.\ Jussieu} 2 (2003), 235--333.
\bibitem[Wal80]{Wal80} J.-L.~Waldspurger, Correspondance de Shimura, \emph{J.\ Math.\ Pures Appl.} 59 (1980), 1--132.
\bibitem[Wal90]{Wal90} J.-L. Waldspurger, D\'emonstration d'une conjecture de dualit\'e de Howe dans le cas $p$-adique, $p\ne2$, in: \emph{Festschrift in honor of I. I. Piatetski-Shapiro}, Israel Math. Conf. Proc. 2, Weizmann, 1990, 267--324.
\bibitem[Wal91]{Wal91} J.-L.~Waldspurger, Correspondances de Shimura et quaternions, \emph{Forum Math.} 3 (1991), 219--307.
\bibitem[Wei73]{Wei73} A.~Weil, \emph{Basic Number Theory}, 2nd ed., Springer, 1973; 3rd ed., 1974.
\bibitem[Wei18]{Wei18} M.~H.~Weissman, $L$-groups and parameters for covering groups, \emph{Ast\'erisque} 398 (2018), 33--186.
\bibitem[Wei64]{Wei64} A. Weil, Sur certains groupes d'op\'erateurs unitaires, \emph{Acta Math.} \textbf{111} (1964), 143--211.
\bibitem[Z23]{Z23} J.~Zou, Local metaplectic correspondence and applications, \emph{Math.\ Z.} 305 (2023), no.~3, Paper No.~43; revised version arXiv:2209.05852v3 (2024).
\bibitem[Zel80]{Zel80} A.~V.~Zelevinsky, Induced representations of reductive $p$-adic groups.\ II.\ On irreducible representations of $\GL(n)$, \emph{Ann.\ Sci.\ \'Ec.\ Norm.\ Sup.} 13 (1980), 165--210.
\end{thebibliography}
